\documentclass[10pt, a4paper]{article}

\usepackage[reqno]{amsmath}
\usepackage{amssymb}
\usepackage{csquotes}

% \RequirePackage[T1]{fontenc}
% \RequirePackage[english,ngerman]{babel}
% \RequirePackage[utf8]{inputenc}

\usepackage{makecell}
\usepackage{vmargin}
\usepackage{multirow}
\usepackage{tabularx}

\usepackage{glmatrix}

\setlength{\unitlength}{1cm}
\setlength{\parindent}{0pt}
\setlength{\parskip}{0.4em}
\setpapersize{A4}
\setmarginsrb{20mm}{15mm}{20mm}{30mm}{0pt}{15mm}{10pt}{0mm}

\begin{document}

% Params %%%%%%%%%%%%%%%%%%%%%%%%%%%%%%%%%%%%%%%%%%%%%%%%%%%%%%%%%%%%%%%%%%%%%%%%%%%%
\newVec{vecA}{0,0,0}
\newVec{vecB}{0,1,0}
\newVec{vecC}{-1,0,1}

Given a triangle defined by the following 3 points: $A=\printVecT{\vecA}$, $B=\printVecT{\vecB}$, and $C=\printVecT{\vecC}$\\

\section{Normal Calculation}

%%%%%%%%%%%%%%%%%%%%%%%%%%%%%%%%%%%%%%%%%%%%%%%%%%%%%%%%%%%%%%%%%%%%%%%%%%%%%%%%%%%%%
Calculate the surface normal of the triangle. Assume that the points $A,B,C$ are given in counter-clockwise order (right-hand rule). Don't forget to normalize the surface normal.

\subVec[vecAB]{\vecB}{\vecA}
\subVec[vecAC]{\vecC}{\vecA}
\crossVec[vecN]{\vecAB}{\vecAC}
\lenVec[vecNormScalar]{\vecN}
\normalizeVec[vecNormed]{\vecN}
%
\begin{align*}
    \vec{n} &= \frac{(B-A)\times(C-A)}{\norm{(B-A)\times(C-A)}}\\
    &=
    \frac{
        \printVec{\vecAB} \times
        \printVec{\vecAC}
    }{
        \norm{\printVecT{\vecAB}\times\printVecT{\vecAC}}
    }
    =
    \frac{
        \printCrossPComp{\vecAB}{\vecAC}
    }{
        \norm{\printVecT{\vecAB}\times\printVecT{\vecAC}}
    }\\
    &=
    \frac{\printVec{\vecN}}{\norm{\printVecT{\vecN}}}
    = \frac{\printVec{\vecN}}{\printScalar{\vecNormScalar}}
    = \printVec[1]{\vecNormed}
\end{align*}

\section{Transformations}

%%%%%%%%%%%%%%%%%%%%%%%%%%%%%%%%%%%%%%%%%%%%%%%%%%%%%%%%%%%%%%%%%%%%%%%%%%%%%%%%%%%%%
\newDiagMat{mirrorMat}{1, -1, 1}

Reflect the triangle across the x/z plane. Provide the reflection matrix (without homogeneous coordinates) and use it to transform the points $A,B,C$.

\transformVec[vecAtransformed]{\mirrorMat}{\vecA}
\transformVec[vecBtransformed]{\mirrorMat}{\vecB}
\transformVec[vecCtransformed]{\mirrorMat}{\vecC}

\begin{align*}
    \mathbf{S} = \printMat{\mirrorMat}
\end{align*}

\begin{align*}
    A' = \mathbf{S}\printVec{\vecA} = \printVec{\vecAtransformed},\quad
    B' = \mathbf{S}\printVec{\vecB} = \printVec{\vecBtransformed},\quad
    C' = \mathbf{S}\printVec{\vecC} = \printVec{\vecCtransformed}
\end{align*}

%%%%%%%%%%%%%%%%%%%%%%%%%%%%%%%%%%%%%%%%%%%%%%%%%%%%%%%%%%%%%%%%%%%%%%%%%%%%%%%%%%%%%
\newVec{myAxis}{1,1,1}
\newScalar{myAngle}{120}

Now we want to rotate $\printScalar{\myAngle}^\circ$ around the axis $\printVecT{\myAxis}$ in counter-clockwise fashion (right-hand rule). Provide the rotation matrix $R$ and use it to transform the points $A,B,C$. (Hint: You can take advantage of the fact that this rotation only swaps the coordinate axes and can therefore dispense with the Rodriguez formula.)

\newRotMat[rotationMatrix]{\myAxis}{\myAngle}
\transformVec[vecAtransformed]{\rotationMatrix}{\vecA}
\transformVec[vecBtransformed]{\rotationMatrix}{\vecB}
\transformVec[vecCtransformed]{\rotationMatrix}{\vecC}

\begin{align*}
    \mathbf{R} = \printMat{\rotationMatrix}
\end{align*}
\begin{align*}
    A' = \mathbf{R}\printVec{\vecA} = \printVec{\vecAtransformed},\quad
    B' = \mathbf{R}\printVec{\vecB} = \printVec{\vecBtransformed},\quad
    C' = \mathbf{R}\printVec{\vecC} = \printVec{\vecCtransformed}
\end{align*}

%%%%%%%%%%%%%%%%%%%%%%%%%%%%%%%%%%%%%%%%%%%%%%%%%%%%%%%%%%%%%%%%%%%%%%%%%%%%
\newScalar{scaleValue}{5}

Provide the matrix that represents a uniform 2D scaling by a factor of \printScalar{\scaleValue}. Use extended coordinates.

\newVec{vecScale}{\scaleValue, \scaleValue}
\newScaleMat{myNewMatrix}{\vecScale}

\[
    \mathbf{S} = \printMat{\myNewMatrix}
\]

%%%%%%%%%%%%%%%%%%%%%%%%%%%%%%%%%%%%%%%%%%%%%%%%%%%%%%%%%%%%%%%%%%%%%%%%%%%%
\newVec{vecTranslation}{2, 1}

Provide the matrix that represents a 2D translation by the vector \printVecT{\vecTranslation}.
\hfill \textbf{(1 Punkt)}

\newTransMat{myTranslationMatrix}{\vecTranslation}
\[
    \mathbf{T} = \printMat{\myTranslationMatrix}
\]

%%%%%%%%%%%%%%%%%%%%%%%%%%%%%%%%%%%%%%%%%%%%%%%%%%%%%%%%%%%%%%%%%%%%%%%%%%%%
\newVec{vecOrigin}{0, 0, 1}

Now apply this 2D translation matrix to the origin point. Provide the calculation and the result.

\newTransMat{myTranslationMatrix}{\vecTranslation}
\transformVec{\myTranslationMatrix}{\vecOrigin}
\[
    O' = \printMat{\myTranslationMatrix}
    \printVec{\vecOrigin} =
    \printVec{\transformVecResult}
\]

%%%%%%%%%%%%%%%%%%%%%%%%%%%%%%%%%%%%%%%%%%%%%%%%%%%%%%%%%%%%%%%%%%%%%%%%%%%%
\newVec{vecPoint}{1, 1, 1}
\newVec{vecTranslation}{2, 1}
\newScalarPi{myAngle}{0.5}

Demonstrate, using an example (with concrete numbers), that for a 2D rotation and a 2D translation, it generally makes a difference whether the rotation is applied to a point first and then the translation, or the translation is applied first and then the rotation.

\newTransMat{myTransMat}{\vecTranslation}
\newRotMatTwoD[myRotMat]{\myAngle}
\mulMat[matTR]{\myTransMat}{\myRotMat}
\mulMat[matRT]{\myRotMat}{\myTransMat}
\transformVec[pointTransformedA]{\matTR}{\vecPoint}
\transformVec[pointTransformedB]{\matRT}{\vecPoint}
\(
    \vec{v} = \printVecT{\vecPoint},\quad
    \vec{t} = \printVecT{\vecTranslation},\quad
    \alpha = \printScalarAsPi{\myAngle}
\)
\begin{align*}
    \mathbf{R}_{\alpha}\mathbf{T}_{t}\vec{v} &= \printMat{\myRotMat}
    \printMat{\myTransMat}
    \printVec{\vecPoint}
    = \printMat{\matRT}\printVec{\vecPoint}
    = \printVec{\pointTransformedA}\\
    \mathbf{T}_{t}\mathbf{R}_{\alpha}\vec{v} &= \printMat{\myTransMat}
    \printMat{\myRotMat}
    \printVec{\vecPoint}
    = \printMat{\matTR}\printVec{\vecPoint}
    = \printVec{\pointTransformedB}
\end{align*}

\section{Look At}

%%%%%%%%%%%%%%%%%%%%%%%%%%%%%%%%%%%%%%%%%%%%%%%%%%%%%%%%%%%%%%%%%%%%%%%%%%%%
\newVec{vecCamera}{1, 2, 3}
\newVec{vecDir}{1, 0, 0}
\newVec{vecUp}{0, 1, 0}

Construct a look-at matrix to render a scene from point \(c = \printVecT{\vecCamera}\). The viewing direction should be chosen so that it points exactly in the positive x-direction. Specify \(A\) or \(A^{-1}\). Indicate whether you are specifying \(A\) or \(A^{-1}\). Provide the matrix with concrete values, not variables. If you lack the necessary information, you may choose them freely. In that case, indicate what you chose.

\crossVec[vecRight]{\vecDir}{\vecUp}
\begin{align*}
    \vec{r} = \vec{d}\times \vec{u}
    = \printVec[1]{\vecDir} \times \printVec[1]{\vecUp}
    = \printVec[1]{\vecRight}
\end{align*}
% Create lookat matrix
\newLookAtMat{\vecRight}{\vecUp}{\vecDir}{\vecCamera}
\[
    \mathbf{A} =
    \begin{bmatrix}
        r_x & r_y & r_z & \trans{-r} c\\
        u_x & u_y & u_z & \trans{-u} c\\
        -d_x & -d_y & -d_z & \trans{d} c\\
        0 & 0 & 0 & 1\\
    \end{bmatrix}
    = \printMat{\lookAtMatResult}
\]

\section{Projection}

%%%%%%%%%%%%%%%%%%%%%%%%%%%%%%%%%%%%%%%%%%%%%%%%%%%%%%%%%%%%%%%%%%%%%%%%%%%%
% OBACHT i) is not automated - j) is but needs i) for setup
\newVec{vecNormal}{0.6,0.8,0}
\newScalar{myDelta}{5}

The projection \(\mathcal{P}(p)\) of a point \(p \in \R^{3}\) using a perspective projection \(\mathcal{P}\) with projection center
at the origin onto an image plane with normal \(n\) and distance \(\delta\) from the origin can be described by \(\mathcal{P}(p) = p \frac{\delta}{\trans{p}n}\).\\
Using homogeneous coordinates, this projection can be written as a matrix-vector product.
Give the corresponding matrix.

\[
    \mathcal{P}(p) = \begin{bmatrix}
        1 & 0 & 0 & 0 \\
        0 & 1 & 0 & 0 \\
        0 & 0 & 1 & 0\\
        \frac{n_x}{\delta} & \frac{n_y}{\delta} & \frac{n_z}{\delta} & 0
    \end{bmatrix}
    \text{oder}
    \begin{bmatrix}
        \delta & 0 & 0 & 0 \\
        0 & \delta & 0 & 0 \\
        0 & 0 & \delta & 0\\
        n_x & n_y & n_z & 0
    \end{bmatrix}
\]

For the specific case \(\vec{n}=\printVecT{\vecNormal}\) and \(\delta = \printScalar{\myDelta}\), this results in the so-called standard projection matrix. Provide this matrix.

\newProjMatGen{generatedProjMat}{\vecNormal}{\myDelta}
\[
    \mathcal{P}_{n, \delta} = \printMat[1]{\generatedProjMat}
\]

%%%%%%%%%%%%%%%%%%%%%%%%%%%%%%%%%%%%%%%%%%%%%%%%%%%%%%%%%%%%%%%%%%%%%%%%%%%%
\newVec{vecPoint}{8, 12, -4}

Using this standard projection matrix, project the point \(p = \printVecT{\vecPoint}\) and express the result as a 3D point \(p' = \trans{(x, y, z)}\). Provide the intermediate steps of your calculation.

\newProjMat{defaultProjMat}
\appendToVec{vecPoint}{\vecPoint}{1}
\transformVec[transformedPoint]{\defaultProjMat}{\vecPoint}
\dehomogenVec[dehomogenPoint]{\transformedPoint}
\[
    p' = \mathbf{P}p =
    \printMat{\defaultProjMat}
    \printVec{\vecPoint}
    = \printVec{\transformedPoint}
    \widehat{=}
    \printVec{\dehomogenPoint}
\]

\section{Lighting \& Shading}
    %%%%%%%%%%%%%%%%%%%%%%%%%%%%%%%%%%%%%%%%%%%%%%%%%%%%%%%%%%%%%%%%%%%%%%%
    % Define Variables here
    %%%%%%%%%%%%%%%%%%%%%%%%%%%%%%%%%%%%%%%%%%%%%%%%%%%%%%%%%%%%%%%%%%%%%%%
    \newVec{vecEye}{0,0,0} % Eyepos
    \newVec{vecX}{0,0,-24} % Point to evaluate
    \newVec{vecXNormal}{0,0,1} % Normal at point
    \newVec{vecLsOnePos}{0,6,-18} % light source 1 position
    \newVec{vecLsOneColor}{0,0.7,0} % light source 1 color
    \newVec{vecLsTwoPos}{-15,0,-9} % light source 2 position
    \newVec{vecLsTwoColor}{0,0,0.4} % light source 2 color
    % Matrixes
    \newDiagMat{matA}{0.15, 0.15, 0.15} % ambient term
    \newDiagMat{matD}{0.5, 0.5, 0.5} % diffuse term
    \newDiagMat{matS}{0.8, 0.8, 0.8} % specular term
    % Shine Term
    \newScalar{scalarShine}{20} % shine term
    %%%%%%%%%%%%%%%%%%%%%%%%%%%%%%%%%%%%%%%%%%%%%%%%%%%%%%%%%%%%%%%%%%%%%%%
    % Calculations
    %%%%%%%%%%%%%%%%%%%%%%%%%%%%%%%%%%%%%%%%%%%%%%%%%%%%%%%%%%%%%%%%%%%%%%%
    % l1, l2, v
    \subVec[vecLsOneDir]{\vecLsOnePos}{\vecX}
    \normalizeVec[vecLsOneNorm]{\vecLsOneDir}
    \subVec[vecLsTwoDir]{\vecLsTwoPos}{\vecX}
    \normalizeVec[vecLsTwoNorm]{\vecLsTwoDir}
    \subVec[vecV]{\vecEye}{\vecX}
    \normalizeVec[vecVNorm]{\vecV}
    % r1
    \transposeMat{\vecXNormal}
    \mulMat[matROneNormal]{\vecXNormal}{\transposeMatResult}
    \mulMatScalar[matROneNormalScalar]{\matROneNormal}{2}
    \transformVec[matROneMul]{\matROneNormalScalar}{\vecLsOneNorm}
    \subVec[vecROneResult]{\matROneMul}{\vecLsOneNorm}
    % r2
    \transposeMat{\vecXNormal}
    \mulMat[matRTwoNormal]{\vecXNormal}{\transposeMatResult}
    \mulMatScalar[matRTwoNormalScalar]{\matRTwoNormal}{2}
    \transformVec[matRTwoMul]{\matROneNormalScalar}{\vecLsTwoNorm}
    \subVec[vecRTwoResult]{\matRTwoMul}{\vecLsTwoNorm}
    % l1*n, l2*n
    \dotVec[scalarPlnOne]{\vecLsOneNorm}{\vecXNormal}
    \maxScalar{0}{\scalarPlnOne}
    \mulMatScalar[matDOneScaled]{\matD}{\maxScalarResult}
    \dotVec[scalarPlnTwo]{\vecLsTwoNorm}{\vecXNormal}
    \maxScalar{0}{\scalarPlnTwo}
    \mulMatScalar[matDTwoScaled]{\matD}{\maxScalarResult}
    % v*r1, v*r2
    \dotVec[scalarPvrOne]{\vecVNorm}{\vecROneResult}
    \maxScalar{0}{\scalarPvrOne}
    \powScalar{\maxScalarResult}{\scalarShine}
    \mulMatScalar[matSOneScaled]{\matS}{\scalarPowResult}
    \dotVec[scalarPvrTwo]{\vecVNorm}{\vecRTwoResult}
    \maxScalar{0}{\scalarPvrTwo}
    \powScalar{\maxScalarResult}{\scalarShine}
    \mulMatScalar[matSTwoScaled]{\matS}{\scalarPowResult}
    % Addition of Ambient, Diffuse and Specular
    \addMat[addMatResultOne]{\matA}{\matDOneScaled, \matSOneScaled}
    \addMat[addMatResultTwo]{\matA}{\matDTwoScaled, \matSTwoScaled}
    % Final multiplications
    \transformVec[matColorOne]{\addMatResultOne}{\vecLsOneColor}
    \transformVec[matColorTwo]{\addMatResultTwo}{\vecLsTwoColor}
    \addVec[vecFinalColor]{\matColorOne}{\matColorTwo}
    %%%%%%%%%%%%%%%%%%%%%%%%%%%%%%%%%%%%%%%%%%%%%%%%%%%%%%%%%%%%%%%%%%%%%%%
    % Printing
    %%%%%%%%%%%%%%%%%%%%%%%%%%%%%%%%%%%%%%%%%%%%%%%%%%%%%%%%%%%%%%%%%%%%%%%

    Given a surface in space, the point \(\printVecT{\vecX}\) lies on this surface. At this point, the surface has the normal \(\printVecT{\vecXNormal}\). Evaluate the Phong illumination model at this point, \emph{without} using the Blinn approximation. The following parameters are given:
    \begin{itemize}
        \item Center of projection at position \(\printVecT{\vecEye}\)
        \item Pointlight 1 at position \(\printVecT{\vecLsOnePos}\)
            with color \(\printVecT{\vecLsOneColor}\)
        \item Pointlight 2 at position \(\printVecT{\vecLsTwoPos}\)
            with color \(\printVecT{\vecLsTwoColor}\)
        \item \(C_a = \printMatDiag{\matA}\)
        \item \(C_d = \printMatDiag{\matD}\)
        \item \(C_s = \printMatDiag{\matS}\)
        \item Shininess \(s = \printScalar{\scalarShine}\)
    \end{itemize}

    \begin{align*}
        L(x,\vec{v}) =
        \sum_{y \in pointlights} \left(C_a + C_d\ \max{\left(0, \trans{\vec{l}}\vec{n}\right)}
            + C_s\ \max{\left(0, \trans{\vec{v}}\vec{r}\right)}^s\right) \cdot c_y
    \end{align*}

    % v
    \begin{align*}
        \vec{v} = \frac{eye - x}{\norm{eye - x}}
        = \frac{
            \printVec{\vecEye} -
            \printVec{\vecX}
        }{
            \norm{\printVec{\vecEye} -
            \printVec{\vecX}}
        }
        = \printVec[1]{\vecVNorm}
    \end{align*}

    % l1, l2
    \[\vec{l} = \frac{y - x}{\norm{y - x}}\]
    \begin{minipage}[t]{0.5\textwidth}
        \begin{align*}
            \vec{l}_1 = \frac{
                \printVec{\vecLsOnePos} -
                \printVec{\vecX}
            }{
                \norm{\printVec{\vecLsOnePos} -
                \printVec{\vecX}}
            }
            = \printVec[1]{\vecLsOneNorm}
        \end{align*}
    \end{minipage}%
    \begin{minipage}[t]{0.5\textwidth}
        \begin{align*}
            \vec{l}_2 = \frac{
                \printVec{\vecLsTwoPos} -
                \printVec{\vecX}
            }{
                \norm{\printVec{\vecLsTwoPos} -
                \printVec{\vecX}}
            }
            = \printVec[1]{\vecLsTwoNorm}
        \end{align*}
    \end{minipage}

    % r1, r2
    \[\vec{r} = 2\vec{n}\trans{\vec{n}}\vec{l}-\vec{l}\]
    \begin{minipage}[t]{0.5\textwidth}
        \begin{align*}
            \vec{r}_1 &= 2 \printVec{\vecXNormal}
            \printVecT{\vecXNormal}
            \printVec[1]{\vecLsOneNorm} -
            \printVec[1]{\vecLsOneNorm}\\
            \vec{r}_1 &= 2 \printMat{\matROneNormal}
            \printVec[1]{\vecLsOneNorm} -
            \printVec[1]{\vecLsOneNorm}\\
            \vec{r}_1 &= \printMat{\matROneNormalScalar}
            \printVec[1]{\vecLsOneNorm} -
            \printVec[1]{\vecLsOneNorm}\\
            \vec{r}_1 &= \printVec[1]{\matROneMul} -
            \printVec[1]{\vecLsOneNorm} =
            \printVec[1]{\vecROneResult}
        \end{align*}
    \end{minipage}%
    \begin{minipage}[t]{0.5\textwidth}
        \begin{align*}
            \vec{r}_2 &= 2 \printVec{\vecXNormal}
            \printVecT{\vecXNormal}
            \printVec[1]{\vecLsTwoNorm} -
            \printVec[1]{\vecLsTwoNorm}\\
            \vec{r}_2 &= 2 \printMat{\matRTwoNormal}
            \printVec[1]{\vecLsTwoNorm} -
            \printVec[1]{\vecLsTwoNorm}\\
            \vec{r}_2 &= \printMat{\matRTwoNormalScalar}
            \printVec[1]{\vecLsTwoNorm} -
            \printVec[1]{\vecLsTwoNorm}\\
            \vec{r}_2 &= \printVec[1]{\matRTwoMul} -
            \printVec[1]{\vecLsTwoNorm} =
            \printVec[1]{\vecRTwoResult}
        \end{align*}
    \end{minipage}

    % substitute
    \begin{align*}
        L(x,\vec{v}) &=
            \left(
                C_a + C_d\ \max{\left(0, \printScalar[1]{\scalarPlnOne}\right)} + C_s\ \max{\left(0, \printScalar[1]{\scalarPvrOne}\right)}^{\printScalar{\scalarShine}}
            \right)
            \cdot \printVec{\vecLsOneColor}\\
        &\quad \quad +
            \left(
                C_a + C_d\ \max{\left(0, \printScalar[1]{\scalarPlnTwo}\right)}
                + C_s\ \max{\left(0, \printScalar[1]{\scalarPvrTwo}\right)}^{\printScalar{\scalarShine}}
            \right)
            \cdot \printVec{\vecLsTwoColor}\\
        &=
        \left(
            \printMat{\matA} +
            \printMat[1]{\matDOneScaled} +
            \printMat{\matSOneScaled}
        \right)
        \cdot \printVec{\vecLsOneColor}\\
        &\quad \quad +
        \left(
            \printMat{\matA} +
            \printMat[1]{\matDTwoScaled} +
            \printMat{\matSTwoScaled}
        \right)
        \cdot \printVec{\vecLsTwoColor}\\
        &=
        \printMat{\addMatResultOne} \cdot \printVec{\vecLsOneColor} +
        \printMat{\addMatResultTwo} \cdot \printVec{\vecLsTwoColor}\\
        & =
        \printVec{\vecFinalColor}
    \end{align*}

\end{document}
